\documentclass[11pt]{article}

\usepackage[margin=1in]{geometry}
\usepackage{amsmath,amssymb,amsthm,mathtools}
\usepackage{graphicx}
\usepackage[hidelinks]{hyperref}

\newtheorem{theorem}{Theorem}
\newtheorem{lemma}{Lemma}
\newtheorem{remark}{Remark}

\newcommand{\R}{\mathbb R}
\newcommand{\HH}{\mathfrak H}
\newcommand{\cC}{\mathcal C}
\newcommand{\cX}{\mathcal X}
\newcommand{\cZ}{\mathcal Z}
\newcommand{\cL}{\mathcal L}
\newcommand{\cG}{\mathcal G}
\newcommand{\ip}[2]{\left\langle #1,#2\right\rangle}
\newcommand{\norm}[1]{\left\lVert #1\right\rVert}

\title{A Higher-Regularity Cochain Extension by Contracting Homotopy}
\author{Automated Math Discovery Candidate Proof}
\date{August 9, 2026}

\begin{document}
\maketitle

\begin{center}
\fbox{\parbox{0.91\textwidth}{\textbf{Status.}
Candidate full solution to Open question 1 of Nilsson--Pitassi,
likely valid pending expert review.  The result answers the literal
fixed-domain boundedness question.  It does not claim a sharp bound uniform
over varying families of Lipschitz domains.}}
\end{center}

\section*{Source Question}

Nilsson and Pitassi ask in Section 6, Open question 1, pages 19--20 of
arXiv:2604.04927v2~\cite{NilssonPitassi}, whether their cochain extension has
a strict higher-regularity analogue.  For $m\in\mathbb N_0$, put
\[
 \cX_m^k(\Omega):=H^{(m,m)}\Lambda^k(\Omega)
 =\{\omega\in H^m\Lambda^k(\Omega):d\omega\in
 H^m\Lambda^{k+1}(\Omega)\}
\]
with its graph norm, and
\[
 \cC_m^k(\Omega):=\cX_m^k(\Omega)\cap\HH^k(\Omega)^\perp,
\]
where $\HH^k(\Omega)$ is the absolute $L^2$ harmonic space used in the
source paper.  The question asks for bounded operators
\[
 E_m^k:\cC_m^k(\Omega)\longrightarrow\cX_m^k(K)
\]
which extend the datum and commute strictly with $d$.  It also asks whether,
when $\Omega\Subset K$, the extensions can have vanishing outer traces and
therefore extend by zero to the whole space.

\begin{center}
\includegraphics[width=\textwidth]{figures/open_problem_crop_page19.png}

\medskip
\includegraphics[width=\textwidth]{figures/open_problem_crop_page20.png}
\end{center}

\section*{Proof Intuition}

The construction separates the problem into an internal homological step on
$\Omega$ and an arbitrary degreewise Sobolev extension.  Maximal regularity
for the de Rham complex implies that the harmonic-orthogonal graph-Sobolev
complex $\cC_m^\bullet(\Omega)$ is exact.  Since it is an exact Hilbert
complex with closed differentials, it has a bounded contracting homotopy
$h$, meaning
\[
                         dh+hd=I.
\]
Let $S^k$ be any bounded extension in degree $k$; it need not commute with
$d$.  The two-term formula
\[
                 E_m^k=dS^{k-1}h^k+S^kh^{k+1}d
\]
repairs it.  Restriction to $\Omega$ gives $dh+hd=I$, while applying $d$ to
the formula makes the cochain identity automatic because $d^2=0$.

The only delicate point is exactness at higher regularity.  An exact $H^m$
form has an $H^{m+1}$ potential by Costabel--McIntosh.  That potential need
not initially be orthogonal to the $L^2$ harmonic fields, but their smooth
cohomology representatives allow a finite-dimensional closed correction
which enforces the orthogonality without changing the derivative.

\section*{Main Result}

\begin{theorem}\label{thm:main}
Let $\Omega\subset\R^n$ be a bounded Lipschitz open set and let
$K\subset\R^n$ be a bounded open set containing $\Omega$.  For every
$m\in\mathbb N_0$ and $0\leq k\leq n$, there is a bounded linear operator
\[
 E_m^k:\cX_m^k(\Omega)\cap\HH^k(\Omega)^\perp
       \longrightarrow \cX_m^k(K)
\]
such that
\begin{align*}
 (E_m^k\omega)|_\Omega&=\omega,\\
 dE_m^k\omega&=E_m^{k+1}(d\omega).
\end{align*}
Here the spaces in degrees outside $\{0,\dots,n\}$ and the corresponding
operators are understood to be zero.

If $\Omega\Subset K$, the operators may be chosen so that every
$E_m^k\omega$ is supported in a fixed compact subset of $K$.  Consequently
all outer traces which occur in the graph-Sobolev setting vanish, and
extension by zero defines a bounded operator into
$H^{(m,m)}\Lambda^k(\R^n)$.
\end{theorem}

\begin{remark}
The theorem is slightly stronger topologically than the question: the
construction does not use any condition on $K\setminus\overline\Omega$.
Its norm is finite for each fixed domain pair.  No quantitative uniformity
over varying domain families is asserted.
\end{remark}

\section*{The Harmonic-Orthogonal Sobolev Complex}

We use two standard facts.  First, the $L^2$ Hodge decomposition in the
notation of the source paper is
\begin{equation}\label{eq:hodge}
 L^2\Lambda^k(\Omega)
 =\HH^k(\Omega)\oplus d(H\Lambda^{k-1}(\Omega))
  \oplus\ker(d|_\Omega)^\perp.
\end{equation}
Second, Costabel and McIntosh prove maximal regularity and regularity of
cohomology on bounded Lipschitz domains~\cite[Corollary 4.7 and
Theorem 4.9]{CostabelMcIntosh}.  In particular:
\begin{enumerate}
\item if $z\in H^m\Lambda^k(\Omega)$ is exact with an $L^2$ potential, then
there is $v\in H^{m+1}\Lambda^{k-1}(\Omega)$ with $dv=z$;
\item there is a finite-dimensional space
$\cG^k\subset C^\infty(\overline\Omega,\Lambda^k)$ of closed forms which
represents de Rham cohomology at every Sobolev order.
\end{enumerate}

\begin{lemma}\label{lem:exact}
For every $m\in\mathbb N_0$, the spaces $\cC_m^k(\Omega)$ form an exact
Hilbert complex under $d$.
\end{lemma}

\begin{proof}
The space $\cC_m^k(\Omega)$ is closed in the Hilbert graph space
$\cX_m^k(\Omega)$.  If $q\in\HH^{k+1}(\Omega)$ and
$\omega\in\cC_m^k(\Omega)$, the Green identity used in the source paper
gives
\[
 \ip{d\omega}{q}_{L^2}
 =\ip{\operatorname t\omega}{\operatorname n q}_{\partial\Omega}
   -\ip{\omega}{\delta q}_{L^2}=0.
\]
Thus $d\omega\perp\HH^{k+1}(\Omega)$, and $\cC_m^\bullet$ is a complex.

Let $z\in\cC_m^k(\Omega)$ satisfy $dz=0$.  By
\eqref{eq:hodge}, the two conditions $dz=0$ and
$z\perp\HH^k(\Omega)$ imply that $z=du_0$ for some
$u_0\in H\Lambda^{k-1}(\Omega)$.  Costabel--McIntosh maximal regularity,
applied with Sobolev indices $s=m$ and $t=0$, yields
$v\in H^{m+1}\Lambda^{k-1}(\Omega)$ with $dv=z$.  In particular
$v\in\cX_m^{k-1}(\Omega)$.

It remains to impose $v\perp\HH^{k-1}(\Omega)$.  Let $P_{\HH}$ be the
$L^2$ orthogonal projection onto $\HH^{k-1}(\Omega)$.  The restriction
\[
                   P_{\HH}|_{\cG^{k-1}}:
                   \cG^{k-1}\longrightarrow\HH^{k-1}(\Omega)
\]
is an isomorphism.  Indeed, both spaces have dimension equal to the relevant
Betti number.  If $g\in\cG^{k-1}$ has $P_{\HH}g=0$, then $g$ is closed and
harmonic-orthogonal, so \eqref{eq:hodge} makes it exact.  Maximal regularity
gives a Sobolev potential, and the direct-sum property of $\cG^{k-1}$ in
\cite[Theorem 4.9]{CostabelMcIntosh} forces $g=0$.

Set
\[
 g=(P_{\HH}|_{\cG^{k-1}})^{-1}P_{\HH}v,
 \qquad u=v-g.
\]
Then $g$ is smooth and closed, hence $u\in\cC_m^{k-1}(\Omega)$, and
$du=dv=z$.  This proves exactness.  At degree zero, a closed function is
locally constant and therefore belongs to $\HH^0(\Omega)$, so the same
statement says that the harmonic-orthogonal cycle space is zero.
\end{proof}

\section*{A Bounded Contracting Homotopy}

\begin{lemma}\label{lem:contraction}
There are bounded linear maps
\[
 h^k:\cC_m^k(\Omega)\longrightarrow\cC_m^{k-1}(\Omega)
\]
such that
\begin{equation}\label{eq:contract}
                     dh^k+h^{k+1}d=I
                     \quad\hbox{on }\cC_m^k(\Omega).
\end{equation}
\end{lemma}

\begin{proof}
Let
\[
 \cZ_m^k=\ker(d:\cC_m^k\to\cC_m^{k+1}),
 \qquad \cL_m^k=(\cZ_m^k)^\perp_{\cC_m^k},
\]
where orthogonality is taken in the graph-space Hilbert norm.  Lemma
\ref{lem:exact} implies that
\[
       d:\cL_m^{k-1}\longrightarrow\cZ_m^k
\]
is a bounded bijection.  Its inverse $Q^k$ is bounded by the bounded inverse
theorem.  Let $P_{\cZ}^k$ be the graph-orthogonal projection onto
$\cZ_m^k$ and define
\[
                         h^k=Q^kP_{\cZ}^k.
\]
If $\omega=z+\ell$ is the decomposition in
$\cZ_m^k\oplus\cL_m^k$, then
\[
                   dh^k\omega=z,
             \qquad h^{k+1}d\omega=\ell.
\]
Their sum is $\omega$, proving \eqref{eq:contract}.  We set
$h^0=h^{n+1}=0$ at the endpoints.
\end{proof}

\section*{Construction of the Extension}

\begin{proof}[Proof of Theorem \ref{thm:main}]
The universal extension theorem of Hiptmair, Li and Zou supplies, for every
degree, a bounded linear section
\[
 S^k:\cX_m^k(\Omega)\longrightarrow
       \cX_m^k(\R^n)
\]
of restriction~\cite{HiptmairLiZou}; this is also the theorem recalled in
the introduction of~\cite{NilssonPitassi}.  Restrict its values to $K$ and
retain the notation $S^k$.  No cochain property of $S^\bullet$ is used.

For $\omega\in\cC_m^k(\Omega)$, define
\begin{equation}\label{eq:E}
 E_m^k\omega
   :=dS^{k-1}h^k\omega+S^kh^{k+1}d\omega.
\end{equation}
The first term is a closed $H^m$ $k$-form and hence belongs to
$\cX_m^k(K)$; the second belongs there by boundedness of $S^k$.
Consequently \eqref{eq:E} defines a bounded linear map in the graph norm.

Let $r$ denote restriction from $K$ to $\Omega$.  Since $rS^j=I$ and
restriction commutes with the distributional exterior derivative,
Lemma~\ref{lem:contraction} gives
\[
 rE_m^k\omega=dh^k\omega+h^{k+1}d\omega=\omega.
\]
Moreover,
\[
 dE_m^k\omega=dS^kh^{k+1}d\omega
 =dS^kh^{k+1}d\omega+S^{k+1}h^{k+2}d^2\omega
 =E_m^{k+1}(d\omega).
\]
This proves the first assertion.

Now suppose $\Omega\Subset K$.  Choose
$\chi\in C_c^\infty(K)$ equal to one on a neighborhood of
$\overline\Omega$ and replace each whole-space section $S^k$ by
$\chi S^k$.  Multiplication by $\chi$ is bounded on every graph space
$\cX_m^k(\R^n)$, and the new map is still a section of restriction.
Formula \eqref{eq:E} now has support contained in the fixed compact set
$\operatorname{supp}\chi\Subset K$.  It therefore vanishes near
$\partial K$, and its extension by zero belongs to
$\cX_m^k(\R^n)$ with the same bound.  This proves the second assertion.
\end{proof}

\section*{Verification and Limitations}

The proof uses no computation.  The identities following
\eqref{eq:E} are formal; the substantive review point is Lemma
\ref{lem:exact}, specifically the compatibility between the absolute
harmonic representatives in~\cite{NilssonPitassi} and the smooth de Rham
cohomology complement in~\cite{CostabelMcIntosh}.  Both represent ordinary
de Rham cohomology on a bounded Lipschitz domain, which makes
$P_{\HH}|_{\cG}$ an isomorphism.

The result gives bounded operators for each fixed $\Omega$, $K$, and $m$.
The abstract bounded inverse used to construct $h$ does not by itself give a
sharp estimate uniform over a varying class of domains.  Thus, if the phrase
``analogue of Theorem 1'' was intended to demand the same quantitative
dependence on a uniform Lipschitz character, that stronger quantitative
problem is not resolved here.

\paragraph{Novelty check.}
On August 9, 2026, bounded searches used the exact source title and close
variants of ``strict higher-regularity analogue,'' ``higher-regularity
cochain extension,'' $H^{(m,m)}$ with ``cochain extension,'' and ``bounded
contracting homotopy Sobolev de Rham Lipschitz extension.''  They found the
three papers cited below but no explicit arbitrary-ambient higher-regularity
cochain section answering Open question 1, and no occurrence of formula
\eqref{eq:E} in this setting.  Novelty confidence is moderate because the
argument is a short synthesis of established analytic and homological
ingredients.

\paragraph{Human-review recommendation.}
Verify Lemma~\ref{lem:exact} first, then confirm that the intended question
does not silently require uniform constants over domain families.  Once
exactness is accepted, Lemma~\ref{lem:contraction} and formula
\eqref{eq:E} complete the fixed-domain question.

\begin{thebibliography}{9}

\bibitem{NilssonPitassi}
E.~Nilsson and S.~Pitassi,
\emph{Uniformly Bounded Cochain Extensions and Uniform Poincar\'e
Inequalities}, arXiv:2604.04927v2, 2026.

\bibitem{CostabelMcIntosh}
M.~Costabel and A.~McIntosh,
\emph{On Bogovski\u{\i} and regularized Poincar\'e integral operators for
de Rham complexes on Lipschitz domains},
Math. Z. \textbf{265} (2010), 297--320; arXiv:0808.2614v2.

\bibitem{HiptmairLiZou}
R.~Hiptmair, J.~Li, and J.~Zou,
\emph{Universal extension for Sobolev spaces of differential forms and
applications}, J. Funct. Anal. \textbf{263} (2012), 364--382.

\end{thebibliography}

\end{document}
